
%     % 1. How can automated model checking and other formal-verification 
%     %     techniques be used to verify different properties of 
%     %     self-replicating robotic systems, and what adaptations 
%     %     can be done for practical implementation?
%
%         In this paper we described usage of NuSMV model checker for 
%         three chosen aspects of SRRS domain:
% \begin{enumerate}
%     
%     \item Checking point and region reachability during local assembly process 
%         and using counterexample to control the movement.
%
%     \item Checking of various robot configurations (block sequences) based on reachability
%     to choose the eligible for the system and environment.
%
%     \item Checking of global liveness of a system based on distribution of material parts
%     in a 2D field
% \end{enumerate}
%
%     % 2. What systematic methodology can translate the behaviors and 
%     %     structural constraints of self-replicating robots into expressive 
%     %     temporal-logic models and specifications?
%
%     We proposed the NuSMV template-base model generation approach to check the 
%     properties mentioned above for chosen type of SRRS. 
%     We demonstrated a validation methodology of the model checking results by
%     simulation using ROS2 and Gazebo simulation.
%     
%
%     % 3. What limitations -- computational, semantic, or 
%     %     modelling-related -- does model checking face in domain of SRRS,
%     %     and how might these limitations can be mitigated?\\
%
%     We explore the computational difference in different model checking tasks,
%     compared aggregated  and sequential configuration checking, point and region based 
%     targets.
%
%
% Future work will focus on incorporating autonomous perception and adaptive 
% control:
% \begin{itemize}
%     \item Exploration of using model checking for different types of robots, 
%         material parts and environment for chosen self-replication system
%         and analysis of other types of self-replication robotic systems.
%     \item Futher analysis of possibilities to control the system on various levels
%         based on counterexamples from model checker.
%     \item Adaptive changing of the system during self-replication
%         through integration of configuration checking into self-replicating process 
%         based on environment conditions.
%     \item Implementation of model checking to collective behavior in SRRS.
% \end{itemize}
% % integrating visual localization of parts, 
% % dynamic path planning. 
% % By doing so, the SRRS framework can 
% % evolve from a manually orchestrated demonstration toward a truly self
% % -replicating robotic organism capable of operating in unstructured environments.
%
%
% In conclusion, this work serves as both a scientific and engineering 
% contribution: 
% scientifically, as a platform for testing hypotheses about 
% self-replicating systems;
% and technically, as a fully functional ROS2 - Gazebo system implementing mechanical 
% assembly. 
% The methods and code we developed--spanning from the launch orchestration
% to the detailed node-level control -collectively provide a reproducible foundation 
% for future research in autonomous manufacturing and self-replication robotics.
%





This work investigated the application of automated model checking and 
to the domain of self-replicating robotic systems.
Self-replicating systems, by their nature, involve dynamic configuration changes, high 
structural and behavioral complexity.
Traditional analysis, simulation or real-based experiments often fails to guarantee the 
correctness of such systems due to the large number of states with logical errors 
or contradictions including interactions between robot, materials and environment. 
The work presented here aimed to fill this gap by employing NuSMV symbolic model checker, 
to rigorously verify selected properties of SRRS under varying structural and 
operational conditions.

The thesis demonstrated how automated model checking can be applied to 
verify three aspects of SRRS behavior. 
Each aspect corresponds to a different level of system abstraction -- 
local, structural, and global, illustrating the generality of the approach.

The first application, described in Chapter \ref{ch:ModelCheckingLocal}, 
focused on verifying point and region reachability during 
the local assembly process of modular robots. By encoding the robot
actuation and attachment logic in NuSMV and defining reachability goals 
as temporal logic specifications, the system could determine whether a 
particular configuration or spatial position was achievable. When a 
specification was violated, NuSMV generated a counterexample trace, 
which then can be used to guide  movement. 
This approach provided not only verification but also a 
foundation for reactive control based on model-checking feedback.

The second verification task addressed the global dynamics of the self
-replication process by modeling the distribution of material parts in 
a 2D workspace was demonstrated in Chapter \ref{ch:global}. 
Here, temporal logic was used to verify whether the 
system preserved liveness -- the ability to eventually complete replication 
given finite resources. The model checking results confirmed that 
certain distribution conditions led to deadlocks, while others 
guaranteed successful replication cycles.

The third aspect of SRRS model checking shown in Chapter \ref{ch:configurationChecking}
and concerned the feasibility of robot configurations.
Using the same modeling framework, the approach evaluated multiple block 
sequences to determine which configurations were valid and functionally reachable 
within the system mechanical and environmental constraints. 
This enabled an automated selection of eligible robot designs -- a task 
traditionally requiring extensive simulation or manual inspection.
The results showed that the NuSMV checker could systematically reject infeasible 

% These  NuSMV-based formal verification approaches can 
% provide structured, interpretable, and repeatable means of analyzing SRRS 
% behavior.
One of the purposes of the research centered on developing a systematic methodology 
to represent SRRS behavior and constraints in a form suitable for temporal-logic analysis.
To achieve this, in Chapter \ref{ch:configurationCheckingPy} 
we introduced a template-based NuSMV model generation 
approach, implemented through Python scripts that automatically translate 
high-level robot specifications into  NuSMV code.

The approach used parameterized templates to encode repeated components for robot
blocks, and to define their logical interactions.
This enabled scalable generation of models representing different robot 
structures without manual redefinition. 
By combining this template-based method with developed temporal-logic specifications
for reachability and liveness, demonstrated in Chapter \ref{ch:reachabilitySpec},
it became possible to evaluate entire families of designs under  
specified verification criteria.

Furthermore, in Chapter \ref{ch:simulation} of this paper we integrated ROS2 
and Gazebo simulations to validate the results of model checking. 
Counterexamples and verified traces were interpreted as executable motion 
sequences and simulated using robot, parts and environment, developed based on the
design principles described in \cite{Jenett2019MaterialRobot}. 

This verification loop -- from formal model to physical simulation --
ensured that verified properties held not only in abstract state space 
but also in realistic dynamic conditions.
This methodological framework provides a reproducible foundation for 
integrating formal verification into the development workflow of modular 
and self-replicating robots.

Model checking suffers from the well-known state explosion problem, 
which becomes bigger in modular robot systems with multiple blocks and 
many degrees of freedom.
Experiments in Chapter \ref{ch:timeAnalysis} showed that the state 
space grows exponentially with each 
added component, but proposed aggregated and sequential configuration checking 
 approach is feasible for moderately 
complex systems, even using personal computer.

% While discrete models allow formal reasoning, they inevitably omit 
% some physical nuances. The methodology mitigated this issue by employing multi-layer 
% modeling, where discrete verification results are validated via 
% continuous simulation in Gazebo. This dual validation helps ensure that 
% verified logical properties correspond to plausible physical outcomes.
%
%
% The NuSMV language, although powerful, presents certain syntactic and 
% expressiveness constraints, especially when modeling hybrid or dynamic topologies.
% The proposed template approach helped mitigate these issues, but future 
% extensions might require adopting tools supporting richer data types or 
% probabilistic reasoning, for example tools like PRISM, or UPPAAL.

Future work in this area can focus on several directions aimed at enhancing both 
the theoretical framework and practical deployment of model-checked SRRS 
systems.
Extending the verification framework to encompass different kinds of 
modular robots, material parts properties, and environmental conditions.
This will allow comparative studies of alternative self-replication paradigms.
Developing higher-level controllers capable of using counterexamples from 
model checking to dynamically guide assembly or repair actions, enabling 
self-correcting replication behaviors.
Incorporating adaptive verification directly into the replication cycle,
allowing the robot to modify its configuration or verification goals 
in response to environmental feedback or detected errors.
Extending the current single-system framework to collective SRRS, 
where multiple robots cooperate in material collection and assembly. 
Formal methods will be applied to verify coordination and 
resource-sharing properties within such distributed systems.


In conclusion, this work makes both scientific and engineering contributions to 
the field of autonomous manufacturing and self-replicating robotics. 
Scientifically, it provides a novel framework for integrating formal verification 
into the analysis of self-replicating processes.
Technically, it delivers a working ROS2--Gazebo implementation of a modular self-
replicating robot whose behavior can be verified, tested, and reproduced.
The developed methods and codebase -- from NuSMV model generation scripts to 
simulation orchestration—constitute a reproducible experimental platform 
for future research in self-replication, modular robotics, and formal verification.
This integration of model checking and robotic simulation paves 
the way toward reliable, verifiably correct systems, 
capable not only of constructing themselves but also of proving the 
correctness of their own actions.
The combination of model checking, simulation validation, 
and adaptive control proposed in this thesis represents a step toward closing 
the gap between theoretical models of self-replication and their practical 
realization in Self-replicating robotic systems.
